\chapter{Esophageal pH Monitoring}

\section*{Overview}
Esophageal pH monitoring measures how often and how long stomach acid or reflux reaches the esophagus. It helps determine whether symptoms are associated with reflux, particularly when the diagnosis is uncertain or symptoms continue despite treatment.

\section*{Why It Is Done}
It helps confirm or characterize GERD when symptoms are unclear, endoscopy is normal, symptoms persist despite treatment, or anti-reflux surgery is being considered.

\section*{How It Is Performed}
A thin catheter passed through the nose or a wireless capsule attached during endoscopy records acid exposure for approximately 24 to 96 hours. The patient logs meals, symptoms, sleep, and medicines.

\section*{Preparation}
The center may request fasting and may give instructions about acid-suppressing medicines. Do not change these medicines unless the ordering clinician provides a plan, because the test question affects preparation.

\section*{Risks and Limitations}
Temporary nose or throat discomfort, gagging, and minor bleeding can occur. Catheter and wireless studies have different limitations, and results depend on the recording period and symptom diary.

\section*{Results and Interpretation}
Results include acid-exposure time and association between reflux episodes and symptoms. Interpretation depends on whether acid-suppressing medicine was continued or stopped and on other findings such as endoscopy.

\section*{Aftercare and Recovery}
Follow instructions for returning the recorder or for removal of a wireless capsule. Contact the team for severe chest pain, persistent vomiting, or difficulty swallowing.

\section*{Contraindications and Precautions}
The team reviews swallowing ability, nasal or esophageal obstruction, recent upper gastrointestinal injury, bleeding risk, and whether acid-suppressing medicines should be continued for the clinical question. Do not change medicines without instructions.

\section*{When to Seek Medical Care}
Follow the procedure team's specific instructions. Seek urgent medical assessment for severe or worsening pain, uncontrolled bleeding, fever or signs of infection, trouble breathing, chest pain, fainting, new weakness or confusion, inability to urinate, or any rapidly worsening symptom. The appropriate warning signs depend on the procedure and the patient's medical condition.

\section*{References}
\begin{enumerate}
  \item MedlinePlus. Esophageal pH Monitoring. U.S. National Library of Medicine; 2025.
  \item Current Medical Diagnosis and Treatment. McGraw-Hill; 2025.
\end{enumerate}
\clearpage
